% Part: computability
% Chapter: computability-theory
% Section: s-m-n

\documentclass[../../../include/open-logic-section]{subfiles}

\begin{document}

\olfileid{cmp}{thy}{smn}
\olsection{$s$-$m$-$n$ উপপাদ্য}

\begin{explain}
পরের উপপাদ্যটি এমন এক কারণে “$s$-$m$-$n$ উপপাদ্য” নামে পরিচিত, যা একটু
পরেই স্পষ্ট হবে। কঠিন কাজটি হলো উপপাদ্যটি ঠিক কী বলছে তা বোঝা; বক্তব্যটি
বুঝে গেলে সেটিকে বেশ স্বাভাবিক বলে মনে হবে।
\end{explain}

\begin{thm}
\ollabel{thm:s-m-n}
  প্রত্যেক জোড়া স্বাভাবিক সংখ্যা $n$ ও~$m$-এর জন্য এমন একটি আদিম
  পুনরাবৃত্ত অপেক্ষক~$s^m_n$ আছে, যাতে প্রত্যেক অনুক্রম
  $e$, $a_0$, \dots, $a_{m-1}$, $y_0$ ,\dots, $y_{n-1}$-এর জন্য
  \[
  \cfind{s^m_n(e, a_0, \dots, a_{m-1})}[n](y_0, \dots, y_{n-1}) \simeq
  \cfind{e}[m+n](a_0, \dots, a_{m-1}, y_0, \dots, y_{n-1}).
\]
\end{thm}

\begin{explain}
$s^m_n$-কে \emph{প্রোগ্রামের} উপর ক্রিয়াশীল বলে ভাবলে সুবিধা হয়। অর্থাৎ,
$s^m_n$ কোনো $(m+n)$-স্থানী অপেক্ষকের প্রোগ্রাম~$e$ এবং আগে থেকে স্থির
নিবেশ $a_0$, \dots, $a_{m-1}$ গ্রহণ করে; তারপর অবশিষ্ট আর্গুমেন্টগুলির
$n$-স্থানী অপেক্ষকের একটি প্রোগ্রাম $s^m_n(e, a_0, \dots, a_{m-1})$
ফেরত দেয়। \iftag{TMs}{$e$-কে কোনো টুরিং যন্ত্রের বর্ণনা হিসেবে ভাবলে,
$s^m_n(e, a_0, \dots, a_{m-1})$ হলো সেই টুরিং যন্ত্র, যা নিবেশ $y_0$,
\dots,~$y_{n-1}$-এর আগে $a_0$, \dots,~$a_{m-1}$ বসিয়ে নিবেশ-স্ট্রিংটি
তৈরি করে এবং~$e$ চালায়। তখন প্রত্যেক $s^m_n$ কেবল উপযুক্ত টুরিং যন্ত্রের
সংকেত খুঁজে দেওয়া একটি আদিম পুনরাবৃত্ত অপেক্ষক।}{}
% BN-SRC-142 TARGET-CORRECTION-BEGIN
\par\smallskip\noindent\textnormal{\small\textbf{উৎস-সংশোধন (BN-SRC-142):} হিমায়িত ইংরেজি ব্যাখ্যায় উপপাদ্যের প্রোগ্রাম-সূচক~$e$-এর জায়গায় দুবার~$x$ এসেছে—একবার $s^m_n(x, a_0, \dots, a_{m-1})$-এ এবং একবার টুরিং যন্ত্রের বর্ণনায়। বাংলা পাঠে উভয় জায়গায়~$e$ রাখা হয়েছে, যাতে বক্তব্যটি উপপাদ্যের সমীকরণের সঙ্গে সঙ্গতিপূর্ণ থাকে।} % BN-SRC-142 TARGET-CORRECTION-END
\end{explain}

\end{document}
